\documentclass[10pt]{article}

% Specify the margins and text size.
\setlength{\textwidth}{6.7in}
\setlength{\textheight}{9.5in}
\setlength{\oddsidemargin}{0pt}
\setlength{\evensidemargin}{0pt}
\setlength{\topmargin}{0pt}
\setlength{\hoffset}{-.1in}
\setlength{\voffset}{-1in}

\setlength{\parskip}{5pt}
\setlength{\parindent}{0pt}

% Load some fonts and symbol packages
\usepackage{latexsym}
\usepackage{pifont}       % contains 'star' symbol for counterinsurgency handbook title
\usepackage{yfonts} 
\usepackage{amsmath}
\usepackage{amsfonts}

\usepackage{graphicx}     % actually, this is loaded by pstricks
\usepackage[T1]{fontenc}
\usepackage{ifthen}
\usepackage{pstricks,pst-grad,pst-text,pst-node,multido,pst-plot,calc,pst-3dplot}
%\usepackage[all]{xy}
%\usepackage{animate}

% The hyperref package inserts links into the pdf.
\definecolor{MyLinkColor}{rgb}{.1,.2,1}
\definecolor{MyCiteColor}{rgb}{.1,1,.2}
\definecolor{MyURLColor}{rgb}{.4,.4,.4}
\usepackage[backref=true,pagebackref=false,hyperindex,colorlinks=true,
  linkcolor=MyLinkColor,urlcolor=MyURLColor]{hyperref}


% The tweaklist package is something I found on the web.  It provides a simple interface
% for making changes to spacing used in the itemize and enumerate environments.  Comment
% this out if you don't care to use tweaklist.
\usepackage{tweaklist}
\renewcommand{\itemhook}{\setlength{\parskip}{2pt}\setlength{\parsep}%
{1pt}\setlength{\topsep}{0pt}\setlength{\itemsep}{0pt}}

\newcommand{\U}{\underline{\hspace{5pt}}}

\usepackage{listings}
\newcommand{\Z}{\hphantom{0}}

\begin{document}
\pagestyle{empty}


\href{http://www.su.edu}{\includegraphics[height=1.75cm]{sulogo.eps}}
\vspace{-1.69cm}

\hfill {\small
\begin{tabular}{cl}
& Math 207\\
& Introduction to Statistics\\
\end{tabular}
}
\setlength{\baselineskip}{1.05\baselineskip}
\bigskip

\begin{center}
\textbf{\large  Exam II Review Exercises (in-class version)}
\end{center}
\medskip

1. \textbf{Standard Normal Curve}:  Sketch and compute the specified area under the normal curve. 
Write down the R command that you use.

\hspace{10pt} a) $z < 2$
\vspace{.4in}

\hspace{10pt} b) $z > 0$
\vspace{.5in}

\hspace{10pt} c) $1 < z < 2$
\vspace{.5in}

\hspace{10pt} d) What $z$ is greater than 99.9\% of all $z$ scores?
\vspace{.5in}

2. \textbf{Heights}:  Suppose that for a particular population
subjects' heights average out to 64 inches with an sd of 6 inches and have
a normal distribution.

\hspace{10pt} a) What percent of the heights were below 72 inches?
\vspace{1.1in}

\hspace{10pt} b) What percent were above 60 inches?
\vspace{1.1in}

\hspace{10pt} c) What percent were between 60 and 72 inches?
\vspace{1.1in}

\hspace{10pt} d) What height was greater than 99.9\% of all the heights?
\vfill
\eject

3. \textbf{Linear Regression}:
In a land animals data set, body weights  average out to 4278 kg with an sd of 16,480 kg.
Brain weights average out to 0.5 kg with an sd of 1.3 kg. The correlation is $-0.005$

\hspace{20pt} a) Write down the equation for the regression line for predicting brain weight
from animal weight.
\vspace{1.0in}

\hspace{20pt} b) Sketch the regression line. 
\vspace{1.25in}

\hspace{20pt} c) 
Estimate the average brain size of an animals that weigh 1 kg (a guinea pig),
465 kg (a cow) and 87,000 kg (brachiosaurus).  The true values are $0.0055$ kg, $0.423$ kg and $0.1545$ kg.
Why does the regression line do such a poor job?
\vspace{2.75in}

\hspace{20pt} d) Estimate the weight of an animal with a brain size of $0.157$ kg (jaguar).
The true value is  100~kg.
\vfill
\eject

4. \textbf{Concept of Correlation}:

\hspace{20pt} a) For a representative sample of people age 25 and older, would the correlation between
age and systolic blood pressure be positive or negative?  Explain.
\vspace{.5in}

\hspace{20pt} b) The correlation between blood pressure and the number of cars that a person has ever owned 
turns out to be positive.
How do you account for this association?
\vspace{.5in}

\hspace{20pt} c)  Suppose that your internet bill is always 30\% of your phone bill.
What can you say about the correlation?
\vspace{.5in}

\hspace{20pt} d)  Suppose that your power bill is always twice as much as your phone bill.
What can you say about the correlation?
\vspace{.5in}

5. \textbf{Regression and the Normal Curve}: In the airquality data set, we find that\vspace{-5pt}
\begin{align*}
\mbox{average temperature}   & \approx \mbox{$78^{\circ}\mbox{F}$,}    & \mbox{sd} & \approx \mbox{$9^{\circ}\mbox{F}$}\\[3pt]
\mbox{average solar radiation}  & \approx \mbox{186\; lang,}  & \mbox{sd} & \approx 90\; \mbox{lang}
   \hspace{15pt} r\approx 0.28
\end{align*}\vspace{-20pt}

\hspace{10pt} a) Predict the wind speed on a day in which the high temperature is 90$^{\circ}\mbox{F}$.
\vspace{.5in}

\hspace{10pt} b) Among 90 degree days, about 95\% had radiation values in what range?
\vspace{1in}

6. \textbf{NHANESIII Data Set}:

\hspace{10pt} a) In the exam data set, what information does the variable \texttt{PEP13E5A} give?
What kind of variable is it?
\vspace{.25in}


\hspace{10pt} b) In the lab data set, what information does the variable \texttt{PLPSI} give?
What kind of variable is it?  What are its units?
\vspace{.4in}

\hspace{10pt} c) In the lab data set, what information does the variable \texttt{HGP} give?
Find the average of all the non-blank values (as you did for other variables in Report~2).

{\color{blue}\begin{verbatim}
lab =read.csv("/Users/ralphwojtowicz/Desktop/math207/data/lab.csv", header=TRUE)
e1 = lab[lab$HGP != 88888, ]
mean(e1$HGP, na.rm=TRUE)
\end{verbatim}}

\vfill
\eject
    
7. \textbf{Computing the Correlation and RMS Error}.
Here is some data:
\begin{center}
{\setlength{\tabcolsep}{40pt}\begin{tabular}{cc|cc|c}
$x$ & $y$ & $z_x$ & $z_y$ & $z_x\,z_y$\\[2pt]\hline
3   & 9 & & &\vphantom{\LARGE Y}\\[20pt]
5   & 5 & & & \\[20pt]
7   & 1 & & &\\[20pt]
\end{tabular}}
\end{center}
\hspace{10pt} a) Complete the table then calculate the correlation.
\vspace{1.5in}

\hspace{10pt} b) Someone uses the line $y=2 - 3\,x$ to approximate the data.
Compute the RMS error for this line.
\vspace{1.5in}

\hspace{10pt} c) What is the RMS error for the regression line?


\vfill
\eject
\end{document}

%%%%%%%%%%%%%%%%%%%%%% HERE

\href{http://www.su.edu}{\includegraphics[height=1.75cm]{sulogo.eps}}
\vspace{-1.69cm}

\hfill {\small
\begin{tabular}{cl}
& Math 207\\
& Introduction to Statistics\\
\end{tabular}
}
\bigskip

\begin{center}
\textbf{\large  Exam II Review Exercises}
\end{center}
\medskip

1. \textbf{Standard Normal Curve}:  Sketch and compute the specified area under the normal curve. 
Write down the R command that you use.

\hspace{10pt} a) $z < 0.5$.
{\color{blue}Using R, \texttt{pnorm(0.5)=0.691 = 69.1\%}.}
\vspace{.4in}

\hspace{10pt} b) $z > 1.5$.
{\color{blue}Using R, \texttt{1 - pnorm(1.5)=0.668 = 66.8\%}.}
\vspace{.5in}

\hspace{10pt} c) $-1.5< z < 0.5$.
{\color{blue}Using R, \texttt{pnorm(0.5) - pnorm(-1.5)=0.625 = 62.5\%}.}
\vspace{.5in}

\hspace{10pt} d) What $z$ is greater than 25\% of all $z$ scores?
{\color{blue}Using R, \texttt{qnorm(0.25)=-0.674}.}
\vspace{.5in}

2. \textbf{Blood Pressure}:  Suppose that for a particular population
systolic blood pressures average out to 120 with an sd of 6 and have
a normal distribution.

\hspace{10pt} a) What percent of the scores were below 123?

{\color{blue} If $x=123$, then $z=(123-120)/6=3/6=0.5$.  As in \#1a, the area is $69.1$\%.}
\vspace{0.8in}

\hspace{10pt} b) What percent were above 129?

{\color{blue} If $x=129$, then $z=(129-120)/6=9/6=1.5$.  As in \#1b, the area is $66.8$\%.}
\vspace{0.8in}

\hspace{10pt} c) What percent were between 111 and 123?

{\color{blue} The $z$ scores are $0.5$ and $(111-120)/6=-9/6=-1.5$.  As in \#1c, the area is $62.5$\%.}
\vspace{0.8in}

\hspace{10pt} d) What blood pressure was greater than 25\% of all the blood pressures?

{\color{blue} As in \#1d, the $z$-score is $-0.674$.  Solve $-0.674 = (x-120)/6$ to get $x=116.0$.}
\vfill
\eject

3. \textbf{Linear Regression}:
In a certain class, midterm scores average out to 70 with an SD of 15, as do 
scores on the final.  The correlation between midterm scores and final scores is 
about~$0.50$.   

\hspace{20pt} a) Write down the equation for the regression line for predicting the final exam score.

{\color{blue} $y-70 = 0.5\,\frac{15}{15}\,(x-70)$ which simplifies to $y-70 = 0.5\,(x-70)$.}
\vspace{0.6in}

\hspace{20pt} b) Sketch the regression line.  {\color{blue} The line goes through the point
  $(x,\,y)=(70,\, 70)$ and has $\mbox{slope}=0.5$.}
\vspace{1.25in}

\hspace{20pt} c) 
Estimate the average final exam scores for students whose midterm scores were
85, 40 and 70. %(10 points)

{\color{blue}
If $x=85$, then $y-70 = 0.5\,(85 - 70)=0.5\times 15 = 7.5$.  So, $y=77.5$.

If $x=40$, then $y-70 = 0.5\,(40-70) = 0.5\times (-30) = -15$.  So, $y=55$.

If $x=70$, then $y-70 = 0.5\,(70-70) = 0$.  So, $y=70$.}
\vspace{2.0in}

\hspace{20pt} d) Estimate the midterm score for a student who scored 90 on the final exam.

{\color{blue}We need a new regression line!  Let $y=\mbox{final exam score}$ and 
$x=\mbox{midterm score}$ as in the answers above.  Our new equation is
\[x - 70 = 0.5\,(y-70)\]
If $y = 90$, then $x-70=0.5\,(90-70) = 0.5\times 20 = 10$.  So, $x=80$.
\medskip

Warning:  If we used the equation from a), substituted $y=90$ then solved for $x$, we would 
get $x=110$ which doesn't make sense.}
\vfill
\eject

4. \textbf{Concept of Correlation}:

\hspace{20pt} a) For a representative sample of people age 25 and older, would the correlation between
education level (number of years completed in school) and income be positive or negative?  Explain.

{\color{blue}It would be positive.  On average, completing more years in school (high school diploma, 
bachelors degree, $\dots$) results in a better-paying job.}


\hspace{20pt} b) The correlation between education level and blood pressure turns out to 
be negative?  How do you account for this association?

{\color{blue} Since higher education leads to a better salary and job, perhaps those with 
higher education also have better health care.  Maybe they have studied nutrition.}


\hspace{20pt} c)  Suppose that in a particular population, fathers were always 3" taller than their sons.
What can you say about the correlation?

{\color{blue} The correlation is equal to $1$.  All the data falls on the line $y=x + 3$ where $x$ is
the son's height, $y$ is the father's height and both are measured in inches.}


\hspace{20pt} d)  Suppose that in a particular population, sisters were always 2" shorter than their brothers.
What can you say about the correlation?

{\color{blue} The correlation is equal to $1$.  All the data falls on the line
$y=x-2$ where $x$ is brother's height and $y$ is sisters height.}


5. \textbf{Regression and the Normal Curve}: A survey of trees in a state park found that:
follows. \vspace{-5pt}
\begin{align*}
\mbox{average girth}         & \approx \mbox{13 inches,}    & \mbox{sd} & \approx \mbox{3 inches}\\[3pt]
\mbox{average height}  & \approx \mbox{76 feet,}  & \mbox{sd} & \approx \mbox{6 feet
   \hspace{15pt}} r\approx 0.5
\end{align*}\vspace{-20pt}

\hspace{10pt} a) Predict the height of a tree that has an eight inch girth. 

{\color{blue}The equation
for the regression line is $\mbox{height} - 76= \mbox{0.5}\frac{6}{3}\,(\mbox{girth} - 13)$ which 
simplifies to\\ $\mbox{height} - 76= \mbox{girth} - 13$ or simply $\mbox{height}=\mbox{girth} + 63$.
So, if $\mbox{girth}=8$ inches, then $\mbox{height}=71$ feet.}
\medskip

\hspace{10pt} b) Among all trees with an eight inch girth, about 95\% had a height in what range?

{\color{blue}The RMS for regression is $\mbox{sd}_y\,\sqrt{1-r^2}=6\,\sqrt{1-(0.5)^2}=6\,\sqrt{0.75}\approx 5.2$.
  So, 95\% of the trees had heights in the range $71\pm 2\times 5.2$ feet
  $=71\pm 10.4$ feet.}


6. \textbf{NHANESIII Data Set}: Use the variable classes from our notes (unordered qualitative, 
 ordered qualitative, discrete, continuous) not the ones from the text.

\hspace{10pt} a) In the exam data set, what information does the variable \texttt{PEPPACE} give?
What kind of variable is it?

{\color{blue} It indicates whether or not the patient had a pacemaker (page 204).  It is an unordered, qualitative 
variable (yes or no) although it is recorded as a number in the data set.}
\smallskip


\hspace{10pt} b) In the lab data set, what information does the variable \texttt{PBPSI} give?
What kind of variable is it?  What are its units?

{\color{blue} It gives information about lead levels in the blood. It is a continuous variable.
The units are micro-moles per liter.}
\medskip

\hspace{10pt} c) In the exam data set, what information does the variable \texttt{BMPARML} give?
Find the average of all the non-blank values (as you did for other variables in Report~2).

{\color{blue}It is upper arm length measured in cm.
\begin{verbatim}
exam =read.csv("/Users/ralphwojtowicz/Desktop/math207/data/exam.csv", header=TRUE)
e1 = exam[exam$BMPARML != 8888, ]
mean(e1$BMPARML, na.rm=TRUE)
\end{verbatim}
The result is 30.3 cm.  See Report 1 for a similar example. Be sure to use the location of
exam.csv on your own computer.  See page 200 of the exam data set
for the error code.}

\vfill
\eject
    
7. \textbf{Computing the Correlation and RMS Error}.
Here is some data:
\begin{center}
{\setlength{\tabcolsep}{40pt}\begin{tabular}{cc|cc|c}
$x$ & $y$ & $z_x$ & $z_y$ & $z_x\,z_y$\\[2pt]\hline
0   & 1 & $-1$ & $-1$& 1\vphantom{\LARGE Y}\\[20pt]
2   & 5 & $0$  & $0$ & 0\\[20pt]
4   & 9 & $1$  & $1$ & 1\\[20pt]
\end{tabular}}
\end{center}
\hspace{10pt} a) Complete the table then calculate the correlation.

{\color{blue}
$\mbox{mean}_x = (0 + 2 + 4)/3 = 6/3 = 2$.  $\mbox{sd}_x=\sqrt{(4 + 0 + 4)/2}= sqrt{8/2}=\sqrt{4}=2$.
If $x=0$, then $z=(0-2)/2 = -1$.  The other values are in the table.
\medskip

$\mbox{mean}_y = (1 + 5 + 9)/3=15/3=5$.  $\mbox{sd}_y=\sqrt{(16 + 0 + 16)/2}=\sqrt{32/2}=\sqrt{16}=4$.
If $y=1$, then $z=(1-5)/4 = -1$.   The other values are in the table.
\medskip

$r=\sqrt{(1 + 0 + 1)/2}=\sqrt{1}=1$. }
\bigskip

\hspace{10pt} b) Someone uses the line $y=3\,x - 1$ to approximate the data.
Compute the RMS error for this line.


{\color{blue}
\begin{center}
\begin{tabular}{cccc}
$x$ & $y$ & predicted $y$ & error\\\hline
0 & 1 & $-1$ & $-2$\\
2 & 5 & $5$  & 0\\
4 & 9 & $11$ & 2\\
\end{tabular}
\end{center}

The RMS error is $\sqrt{(4 + 0 + 4)/3}=\sqrt{8/3}\approx 1.6$}
\bigskip

\hspace{10pt} c) What is the RMS error for the regression line?

{\color{blue} It is zero since the points fall on a line.  
We can use the formula:  $\mbox{sd}_y\sqrt{1-r^2} = \mbox{sd}_y\,\sqrt{1-1} = 0$.}
\vfill
\eject

\end{document}
